\chapter{Сборка тестового приложения и базовое профилирование}
\label{ch:baseline}

\section{Параметры исследования}

В качестве тестовой нагрузки используется вычислительный цикл:
$$
y_i = a x_i + y_i.
$$

Программа позволяет изменять количество элементов \texttt{N}, число повторов \texttt{repeats}, шаг доступа \texttt{stride} и тип данных.

\begin{commands}
g++ daxpy.cpp -O3 -std=c++17 -o daxpy
./daxpy 1000000 100 1
/usr/bin/time -l ./daxpy 1000000 100 1
\end{commands}

Здесь \texttt{1000000} --- количество элементов, \texttt{100} --- число повторов, \texttt{1} --- \texttt{stride}.

$$
\mathit{IPC} = \frac{\mathit{Instructions}}{\mathit{Cycles}}.
$$

\section{Измеренные показатели}

\begin{table}[H]
\caption{Базовые показатели (\texttt{./daxpy 1000000 100 1})}
\label{tab:baseline}
\begin{tabularx}{\textwidth}{|X|r|}
	\hline
	Показатель & Значение \\
	\hline
	Время выполнения, с & \texttt{0.04} \\
	\hline
	\texttt{cycles} & \texttt{125609823} \\
	\hline
	\texttt{instructions} & \texttt{721017564} \\
	\hline
	IPC & \texttt{5.7401} \\
	\hline
\end{tabularx}
\end{table}

\section{Вывод}

На втором этапе получены корректно собранная и запускаемая программа, а также набор базовых метрик: время выполнения, cycles, instructions и рассчитанное значение IPC.
